% ----- Consignes exo 8 ----- %
\begin{td-exo}[Sur le problème du sac à dos]\, % 8 
	\begin{enumerate}
		\item Rappeler la formulation classique du problème du sac à dos.

		\item Rajouter des coupes pour améliorer la qualité de la formulation.

		\item Comparer les deux formulations précédentes (relaxations linéaires et solutions exactes) pour les instances données à l'adresse:
		\url{https://artemisa.unicauca.edu.co/johuyortega/instances_01_KP/}
		ou
		\url{https://github.com/JorikJooken/knapsackProblemInstances/}.
	\end{enumerate}
\end{td-exo}

% ----- Solutions exo 8 ----- %
\iftoggle{showsolutions}{ 
	\begin{td-sol}[]\ % 8
		\begin{enumerate}
			\item La formulation classique du problème du sac à dos est la suivante:
			\begin{equation*}
				\begin{cases}
					\max \sum_{i=1}^n p_i z_i \\
					\sum_{i=1}^n w_i z_i \leq W \\
					z_i \in \{0,1\} \quad \forall i = 1, \ldots, n
				\end{cases}
			\end{equation*}
			où \(p_i\) est le profit de l'objet \(i\), \(w_i\) est le poids de l'objet \(i\), \(W\) est la capacité du sac à dos, et \(z_i\) est une variable binaire indiquant si l'objet \(i\) est inclus dans le sac à dos ou non.

			\item Pour améliorer la qualité de la formulation, on peut ajouter des coupes de Gomory ou d'autres types de coupes valides pour le problème du sac à dos. Par exemple, on peut ajouter les coupes suivantes:
			\begin{equation*}
				\sum_{i \in S} z_i \leq |S| - 1 \quad \forall S \subseteq \{1, \ldots, n\} \text{ tel que } \sum_{i \in S} w_i > W
			\end{equation*}
			ces coupes interdisent les combinaisons d'objets qui dépassent la capacité du sac à dos.
			Plus précisément, si un ensemble d'objets \(S\) a un poids total supérieur à \(W\), alors au moins un de ces objets doit être exclu du sac à dos, ce qui est exprimé par la contrainte ci-dessus.
			Une autre coupe possible est la suivante:
			\begin{equation*}
				\sum_{i=1}^n \frac{w_i}{W} z_i \leq 1
			\end{equation*}
			qui est une coupe de couverture linéaire.
			Plus généralement, on peut ajouter des coupes de type ``cover'' ou ``knapsack cover'' qui sont basées sur les propriétés combinatoires du problème du sac à dos.
		\end{enumerate}
	\end{td-sol}
}{}
